\section{总结与延伸}

本讲把 test-time compute 从“多想一会儿”具体化为两个可执行系统。

\begin{enumerate}
    \item \textbf{AlphaCode}：用大规模并行采样覆盖程序空间，用样例、行为聚类和隐藏测试压缩候选；
    \item \textbf{AlphaCode 2}：用更强 foundation model、多个 solver family 与 learned scorer 提高样本效率和提交质量；
    \item \textbf{Search-o1}：在 reasoning 内识别知识缺口，多轮检索并压缩文档，让新证据进入后续推理；
    \item \textbf{Search-R1}：把“何时搜、搜什么、何时停止”进一步视为可通过 RL 学习的策略。
\end{enumerate}

跨越代码搜索与知识搜索后，可以提炼出同一套 agent 设计准则：

\begin{itemize}
    \item 生成器负责扩大候选覆盖，但必须维持语义多样性；
    \item 选择器要基于任务相关的可验证信号，而非生成概率；
    \item 计算预算应按难度和当前不确定性动态分配；
    \item 长轨迹必须保存结构化中间状态，支持局部验证和回溯；
    \item 外部工具会引入新错误面，检索结果、测试与 reward model 自身也要被验证。
\end{itemize}

下一讲转向 agentic evaluation：当任务需要几十分钟到数小时、输出是 CAD、报告、表格或文献综述时，如何测量可靠性、经济价值和真实工作差距。

\end{document}
